% ch3_a §3.1–3.2 (书页 77–82 / p091–p096)
% 第3章 Electron States
\chapter{电子态}

\epigraphbook{谱写部落歌谣的方法有六十九种，\\而其中每一种都无一不是对的。}{KIPLING}

\section{自由电子}

考察固体 Na。每个原子含有 11 个电子，但其中 10 个电子处于紧紧束缚在原子核上的状态，构成一个净正电荷为 $+|e|$ 的离子。在自由原子中，剩下的那个电子在围绕此离子的一个轨道（orbital）上运动。当诸原子聚集成固体时，这些轨道相互交叠并发生相互作用。可以论证，这种交叠是如此广泛，以致于把每个电子都定域在各自原子上的量子化方案必然失效，必须代之以另一种方案：每个电子的波函数是它在所有离子的势中运动的薛定谔方程（Schrödinger equation）的解。于是，我们一开始就区分两类电子：一类是芯电子（core electrons），它们被处理成几乎完全定域化的；另一类是价电子（valence electrons）或传导电子（conduction electrons），假定它们进入 Bloch 态——即延展于整个晶体之中的状态。

在本章中我们将采用单电子模型（one-electron model），即忽略价电子之间因库仑排斥而生的任何相互作用。这些关联（correlation）与交换（exchange）效应将在第 5 章中讨论。芯内部各电子之间、或芯电子与价电子之间的这类效应，通常被当作每个价电子在晶体中运动时所感受到的势 $\mathcal{V}(\mathbf{r})$ 的一部分来处理。于是，$\mathcal{V}(\mathbf{r})$ 被认为是按离子的 Hartree 势或 Hartree--Fock 势计算出来的。

如果假定 $\mathcal{V}(\mathbf{r})$ 是一个常数，就可以做一个极大的简化。这时电子便自动落入平面波态（plane-wave states）
\begin{equation*}
|{\bf k}\rangle=\psi_{\bf k}=e^{i{\bf k}\cdot{\bf r}},
\tag{3.1}
\end{equation*}
其能量为
\begin{equation*}
\mathcal{E}({\bf k})=\frac{\hbar^{2}k^{2}}{2m}.
\tag{3.2}
\end{equation*}

正如在 §1.5 中所见，这些态满足 Bloch 定理。等能面是 ${\bf k}$ 空间中的球面。${\bf k}$ 的允许值在此空间中以 $V/8\pi^{3}$ 的密度分布。但对每一个 ${\bf k}$ 值，实际上有两个电子态，自旋相反。如果实空间中单位体积内有 $ZN$ 个电子（即每个原子 $Z$ 个电子，每个晶胞 1 个原子，单位体积内有 $N$ 个晶胞），那么为了满足 Pauli 原理，我们需要把态一直填充到波数 $k_{F}$，$k_{F}$ 满足
\begin{align*}
\frac{4}{3}\pi k_{F}^{3}\,\frac{2}{8\pi^{3}}&=ZN,\\
\text{即}\quad k_{F}&=(3\pi^{2}ZN)^{1/3}.
\tag{3.3}
\end{align*}

%FIG{36}: {自由电子的费米面。}

我们说费米面（Fermi surface）是一个半径为 $k_{F}$ 的球面。费米能（Fermi energy）或费米能级（Fermi level）是相应于这一分布顶端的能量，
\begin{equation*}
\mathcal{E}_{F}=\frac{\hbar^{2}k_{F}^{2}}{2m}.
\tag{3.4}
\end{equation*}

我们注意到，费米球的半径与 Debye 球的半径相当。把式 (3.3) 与式 (2.53) 相比较，得
\begin{equation*}
k_{F}=\left(\frac{Z}{2}\right)^{1/3}q_{D}.
\tag{3.5}
\end{equation*}

这意味着费米能级上电子的波长与原子间距相当（参看式 (2.55)）：
\begin{equation*}
\lambda_{F}=\frac{2\pi}{k_{F}}\sim\left(\frac{2}{Z}\right)^{1/3}2.6\,r_{s},
\tag{3.6}
\end{equation*}
其中 $r_{s}$ 是 Wigner--Seitz 球的半径。

\section{价电子的衍射}

我们刚刚证明的结果使自由电子模型（free-electron model）露出了破绽：费米能级上的电子所处的态，其波长与晶格间距相当。因此必定存在强的衍射效应，即使 $\mathcal{V}(\mathbf{r})$ 与常数势没有非常显著的差别。让我们把这个势当作作用在自由电子态上的微扰来处理。令
\begin{equation*}
\mathcal{E}({\bf k})=\mathcal{E}_{\bf k}^{0}+\langle{\bf k}|\mathcal{V}|{\bf k}\rangle+\sum_{{\bf k}'}\frac{|\langle{\bf k}|\mathcal{V}|{\bf k}'\rangle|^{2}}{\mathcal{E}_{\bf k}^{0}-\mathcal{E}_{{\bf k}'}^{0}},
\tag{3.7}
\end{equation*}
其中
\begin{equation*}
\mathcal{E}_{\bf k}^{0}\equiv\hbar^{2}k^{2}/2m,
\tag{3.8}
\end{equation*}
并且对势在未微扰态 (3.1) 之间的矩阵元采用了 Dirac 记号。

但 $\mathcal{V}(\mathbf{r})$ 必须具有晶格的周期性。如 §2.6 中所证明，除非 $({\bf k}-{\bf k}')$ 是一个倒格矢（reciprocal lattice vector），否则它的矩阵元为零。由式 (2.78) 立即得到
\begin{equation*}
\mathcal{E}({\bf k})=\mathcal{E}_{\bf k}^{0}+\mathcal{V}_{0}+\sum_{{\bf g}\neq{\bf 0}}\frac{|\mathcal{V}_{\bf g}|^{2}}{\mathcal{E}_{\bf k}^{0}-\mathcal{E}_{\bf k-{\bf g}}^{0}},
\tag{3.9}
\end{equation*}
其中 $\mathcal{V}_{\bf g}$ 是 $\mathcal{V}$ 对应于倒格矢 ${\bf g}$ 的 Fourier 分量（如式 (1.15)）。

问题在于——这样一个展开究竟有多好？有两个条件必须满足：(i) 当 ${\bf g}$ 增大时，Fourier 分量 $\mathcal{V}_{\bf g}$ 应当迅速趋于零；(ii) 在被微扰所混合的诸未微扰态之间，不得存在如下形式的简并
\begin{equation*}
\mathcal{E}_{\bf k}^{0}=\mathcal{E}_{\bf k-{\bf g}}^{0}
\tag{3.10}
\end{equation*}
目前，先把条件 (i) 搁置，留待以后研究。

简并条件 (3.10) 等价于写成
\begin{equation*}
|{\bf k}|=|{\bf k}-{\bf g}|;
\tag{3.11}
\end{equation*}
从几何上看，这意味着 ${\bf k}$ 位于倒格矢 ${\bf g}$ 的垂直平分面上。也就是说——\emph{当 ${\bf k}$ 位于（或接近）区边界时，简单的微扰展开式 (3.9) 便不再有效}。这个公式只能在未微扰态的球面未触及区的边界时使用——实际上，即使对单价金属，它也很少有效。

%FIG{37}: {}

注意，图 37 正是波矢为 ${\bf k}$ 的电子相干 Bragg 衍射到方向 ${\bf k}'={\bf k}-{\bf g}$ 的 Ewald 作图（图 29(b)）。在每一种情形下，几何条件都是：$OPQ$ 应构成一个等腰三角形，且以倒格矢 ${\bf g}$ 为底边 $PQ$。于是我们从代数上证实：晶体内的价电子受到来自晶格的衍射作用，恰好像它们是从晶体外部入射进来的一样。

为了求出区边界邻域内的解，我们必须显式地考察微扰方程。晶格的周期性意味着势只能混合相差一个倒格矢 ${\bf g}$ 的波，因此我们是假定波函数可以展开成级数
\begin{equation*}
\psi_{\bf k}=\sum_{\bf g}\alpha_{\bf k-{\bf g}}e^{i({\bf k}-{\bf g})\cdot{\bf r}}.
\tag{3.12}
\end{equation*}
若把它代入薛定谔方程
\begin{equation*}
\left\{-\frac{\hbar^{2}}{2m}\nabla^{2}+\mathcal{V}(\mathbf{r})\right\}\psi_{\bf k}=\mathcal{E}({\bf k})\psi_{\bf k},
\tag{3.13}
\end{equation*}
并用展开式中的某一项去乘遍整个方程，便得到系数 $\alpha$ 的线性方程组：
\begin{equation*}
\{\mathcal{E}_{\bf k-{\bf g}}^{0}-\mathcal{E}({\bf k})\}\alpha_{\bf k-{\bf g}}+\sum_{{\bf g}'}\mathcal{V}_{{\bf g}-{\bf g}'}\alpha_{\bf k-{\bf g}'}=0
\tag{3.14}
\end{equation*}
其中 ${\bf g}$ 取一切值（包括 ${\bf g}={\bf 0}$ 和 ${\bf g}'={\bf 0}$）。众所周知，只须假定除 $\alpha_{\bf k}$ 之外所有 $\alpha_{\bf k-{\bf g}}$ 都很小，便可由此得到简单的微扰公式 (3.9)。实际上，这时我们取
\begin{equation*}
\alpha_{\bf k-{\bf g}}\approx\frac{\mathcal{V}_{-{\bf g}}}{\mathcal{E}_{\bf k}^{0}-\mathcal{E}_{\bf k-{\bf g}}^{0}}.
\tag{3.15}
\end{equation*}

当 ${\bf k}$ 位于某条区边界附近时——比如说平分矢量 ${\bf G}$ 的那条边界——恰恰就是这个假定失效了。波 $\exp\{i({\bf k}-{\bf G})\cdot{\bf r}\}$ 的系数变得与原来未微扰态的系数同等重要。实际上，此时电子已被晶格作了 Bragg 反射，就好像它是一束从外部入射的电子束一样。条件 (3.11) 等价于 §2.6 中的式 (2.79) 和 (2.80)。为了构成一个定态，我们必须把入射束和衍射束放在同等的地位上处理。

作为近似，让我们忽略 (3.14) 中除这两个系数之外的所有系数。把能量原点移动 $\mathcal{V}_{0}$ 之后，得到
\begin{equation*}
\left.\begin{aligned}
\{\mathcal{E}_{\bf k}^{0}-\mathcal{E}({\bf k})\}\alpha_{\bf k}+\mathcal{V}_{\bf G}\alpha_{\bf k-{\bf G}}&=0,\\
\mathcal{V}_{-{\bf G}}\alpha_{\bf k}+\{\mathcal{E}_{\bf k-{\bf G}}^{0}-\mathcal{E}({\bf k})\}\alpha_{\bf k-{\bf G}}&=0.
\end{aligned}\right\}
\tag{3.16}
\end{equation*}

方程组的行列式是 $\mathcal{E}({\bf k})$ 的一个二次式，其解为
\begin{equation*}
\mathcal{E}^{\pm}({\bf k})=\tfrac{1}{2}(\mathcal{E}_{\bf k}^{0}+\mathcal{E}_{\bf k-{\bf G}}^{0})\pm\tfrac{1}{2}\sqrt{(\mathcal{E}_{\bf k}^{0}-\mathcal{E}_{\bf k-{\bf G}}^{0})^{2}+4|\mathcal{V}_{\bf G}|^{2}}
\tag{3.17}
\end{equation*}
（注意到 $\mathcal{V}_{-{\bf G}}=\mathcal{V}_{\bf G}^{*}$）。这样，态 $e^{i{\bf k}\cdot{\bf r}}$ 与 $e^{i({\bf k}-{\bf G})\cdot{\bf r}}$ 组合成另外两个态 $\psi^{+}$ 和 $\psi^{-}$，其能量分别为 $\mathcal{E}^{+}$ 和 $\mathcal{E}^{-}$。

为简单起见，让我们考虑在一维情形下会发生什么。在 $k=0$ 附近，两个未微扰能量之差大得很，以致我们可以忽略 $4|\mathcal{V}_{\bf G}|^{2}$。我们得到
\begin{equation*}
\mathcal{E}({\bf k})\sim\mathcal{E}_{\bf k}^{0},
\tag{3.18}
\end{equation*}
因而 $\psi^{-}$ 沿着自由电子抛物线。在 $k=\frac{1}{2}G$ 处，也就是区边界处，两个未微扰波具有相同的能量：于是
\begin{equation*}
\mathcal{E}^{-}(\tfrac{1}{2}G)=\mathcal{E}_{\frac{1}{2}G}^{0}-|\mathcal{V}_{\bf G}|.
\tag{3.19}
\end{equation*}
（译注：原书此式右端误印为“$+$”，按式 (3.17) 及上下文“$\psi^{-}$ 的能量被压低”应为“$-$”。）态 $\psi^{-}$ 的能量比自由电子抛物线低了一个量 $|\mathcal{V}_{\bf G}|$。我们在图 38 中的 $A$ 点看到这一点。

%FIG{38}: {一维情形的电子能量：简约区图式。}

在区边界处，态 $\psi^{+}$ 具有能量
\begin{equation*}
\mathcal{E}^{+}(\tfrac{1}{2}G)=\mathcal{E}_{\frac{1}{2}G}^{0}+|\mathcal{V}_{\bf G}|,
\tag{3.20}
\end{equation*}
因而它的能量被抬高到自由电子抛物线之上。若把 $k$ 限制在区之内，我们发现随着 $k$ 的减小，这个态的能量趋于 $\mathcal{E}_{\bf k-{\bf G}}^{0}$。这样，区内的能量便落入两个\emph{能带}（band），中间隔着宽度为 $2|\mathcal{V}_{\bf G}|$ 的带隙（gap）$AB$。

把图 38 与 §1.5 中的图 14 加以比较是有启发性的。那个颇为人为的自由电子简约区图式（reduced zone scheme），已经被周期性晶格势分裂成若干个能带。我们也可以用另一种方式来表示这个结果：把 $\psi^{+}$ 这一支画到第二个区中去——即采用扩展区图式（extended zone scheme）（图 39）。这样更清楚地表明，各个能量是如何由自由电子抛物线经微扰而得出的——只是在区边界处引入了不连续性。这同我们在 §2.2 的图 19(c) 中观察到的现象一样，在那里，两种不同质量的引入使晶格的振动谱发生了分裂。

%FIG{39}: {一维情形的电子能量：扩展区图式。}

为了理解所发生的事情，让我们计算 $k=\frac{1}{2}G$ 处的波函数。假定 $\mathcal{V}_{\bf G}$ 为负。由式 (3.12) 和 (3.16) 我们得到
\begin{align*}
\psi^{-}&=(e^{i\frac{1}{2}Gr}+e^{-i\frac{1}{2}Gr})/\sqrt{2}\\
&=\sqrt{2}\cos\tfrac{1}{2}Gr
\tag{3.21}
\end{align*}
以及
\begin{equation*}
\psi^{+}=\sqrt{2}\,i\sin\tfrac{1}{2}Gr.
\tag{3.22}
\end{equation*}

因此，$|\psi^{-}|^{2}$ 是这样一个函数：它在 $r=0$ 附近很大，而且也在
